\documentclass[12pt]{article}
\usepackage[utf8]{inputenc}
\usepackage[T1]{fontenc}
\usepackage{hyperref}
\title{Yagwoia Language Resources for K-12 Curriculum Development}
\author{Compiled by Claude Sonnet 5.5 (claude-sonnet-5-5) for Andr\'as Kornai}
\date{}
\begin{document}
\maketitle

\section{Introduction}
Yagwoia (also Yeghuye, Kokwaiyakwa; ISO 639-3: ygw) is an Angan language of Papua New Guinea with about 10,000 speakers. It is spoken in the Eastern Highlands, Gulf and Morobe provinces. The search found no dedicated grammar or dictionary and no school materials. The resources below are limited.

\section{Primary Educational Resources}
No primers, textbooks or school programmes were found in search results.

\section{Linguistic and Reference Materials}
\begin{itemize}
\item Wikipedia overview, with a phonology section and dialect names: \url{https://en.wikipedia.org/wiki/Yagwoia_language}
\item Lloyd (1973), ``The Angan Language Family'', which includes Yagwoia word lists and is hosted in the Trans New Guinea database: \url{https://transnewguinea.org/source/lloyd-1973}
\end{itemize}

\section{Religious and Community Resources}
\begin{itemize}
\item Joshua Project profile. It lists the Bible translation status as New Testament: \url{https://joshuaproject.net/people_groups/15943}
\end{itemize}

\section{Digital and Audio Resources}
\begin{itemize}
\item Global Recordings Network page with evangelistic audio and video, including a 40-section Bible lesson set and an MP3 zip: \url{https://globalrecordings.net/en/language/ygw}
\item Wikitongues page for Yagwoia: \url{https://wikitongues.org/languages/ygw/}
\end{itemize}

\section{Curriculum Development Context}
The language is reported as stable. There is a New Testament and recorded audio, and no evidence of school use. A curriculum would need new work on orthography, primers and a dictionary with community speakers.
\end{document}